\documentclass{article}
\usepackage{amsmath}
\usepackage{geometry}
\usepackage{array}

\geometry{a4paper, margin=1in}

\title{Q1 - Mamdani Fuzzy Controller Report}
\author{Producer: MohammadAmin HorriFarahani \\ Student Number: [40216673]}
\date{\today}

\begin{document}

\maketitle

\section{Problem Statement}
We are given a pole-balancing controller with the following crisp inputs:
\begin{itemize}
    \item Pole angle $\theta = -1.5$
    \item Angular velocity $\dot{\theta} = 2$
\end{itemize}
We must estimate the crisp control force $F^*$ using a Mamdani fuzzy inference system with Min-Max logic and Centroid defuzzification over discrete samples.

\section{Step 1: Fuzzification}
We compute the membership values for the inputs based on the piecewise-linear functions provided.

\subsection{Input 1: $\theta = -1.5$}
\begin{itemize}
    \item $\mu_{\text{NEG}}(-1.5) = \frac{0 - (-1.5)}{2} = 0.75$
    \item $\mu_{\text{ZERO}}(-1.5) = 1 - \frac{|-1.5|}{2} = 0.25$
    \item $\mu_{\text{POS}}(-1.5) = 0$
\end{itemize}

\subsection{Input 2: $\dot{\theta} = 2$}
\begin{itemize}
    \item $\mu_{\text{NEG}}(2) = 0$
    \item $\mu_{\text{ZERO}}(2) = 1 - \frac{|2|}{5} = 0.6$
    \item $\mu_{\text{POS}}(2) = \frac{2}{5} = 0.4$
\end{itemize}

\section{Step 2: Rule Evaluation}
We determine the firing strength $\alpha$ for each active rule using the MIN operator. The active antecedent terms are $\theta \in \{\text{NEG}, \text{ZERO}\}$ and $\dot{\theta} \in \{\text{ZERO}, \text{POS}\}$.

\begin{enumerate}
    \item \textbf{Rule 1:} IF $\theta$ is NEG AND $\dot{\theta}$ is ZERO THEN F is N.
    \[ \alpha_1 = \min(0.75, 0.6) = 0.6 \]
    
    \item \textbf{Rule 2:} IF $\theta$ is NEG AND $\dot{\theta}$ is POS THEN F is Z.
    \[ \alpha_2 = \min(0.75, 0.4) = 0.4 \]
    
    \item \textbf{Rule 3:} IF $\theta$ is ZERO AND $\dot{\theta}$ is ZERO THEN F is Z.
    \[ \alpha_3 = \min(0.25, 0.6) = 0.25 \]
    
    \item \textbf{Rule 4:} IF $\theta$ is ZERO AND $\dot{\theta}$ is POS THEN F is P.
    \[ \alpha_4 = \min(0.25, 0.4) = 0.25 \]
\end{enumerate}

\section{Step 3: Implication and Aggregation}
We apply Min-clipping to the consequents and aggregate using Max.

\begin{itemize}
    \item Term \textbf{N}: Clipped at $\alpha = 0.6$.
    \item Term \textbf{Z}: Clipped at $\max(0.4, 0.25) = 0.4$.
    \item Term \textbf{P}: Clipped at $\alpha = 0.25$.
\end{itemize}

\[ \mu_{\text{agg}}(F) = \max \left( \min(\mu_N(F), 0.6), \min(\mu_Z(F), 0.4), \min(\mu_P(F), 0.25) \right) \]

\section{Step 4: Defuzzification}
We use the Riemann sum method over the sample points $F_i \in \{-16, -12, -8, -4, 0, 4, 8, 12, 16\}$.

\begin{center}
\begin{tabular}{|c|c|c|c|}
\hline
$F_i$ & Calculations ($\max(\text{clipping})$) & $\mu_{\text{agg}}(F_i)$ & $\mu_{\text{agg}}(F_i) \cdot F_i$ \\
\hline
$-16$ & $\max(0, 0, 0)$ & 0 & 0 \\
$-12$ & $\max(0.5, 0, 0)$ & 0.5 & $-6.0$ \\
$-8$ & $\max(0.6, 0, 0)$ & 0.6 & $-4.8$ \\
$-4$ & $\max(0.5, 0.4, 0)$ & 0.5 & $-2.0$ \\
$0$ & $\max(0, 0.4, 0)$ & 0.4 & 0 \\
$4$ & $\max(0, 0.4, 0.25)$ & 0.4 & $1.6$ \\
$8$ & $\max(0, 0, 0.25)$ & 0.25 & $2.0$ \\
$12$ & $\max(0, 0, 0.25)$ & 0.25 & $3.0$ \\
$16$ & $\max(0, 0, 0)$ & 0 & 0 \\
\hline
\textbf{Total} & & $\sum = 2.9$ & $\sum = -6.2$ \\
\hline
\end{tabular}
\end{center}

\[ F^* \approx \frac{-6.2}{2.9} \approx -2.1379 \]

\section{Conclusion}
The estimated crisp control force is \textbf{-2.14}.

\section{Problem Statement}
We are tasked with comparing the crisp output of a Zero-Order Sugeno system and a Mamdani system given specific firing strengths and consequent parameters.

\textbf{Given Firing Strengths:}
\[ w_1 = 0.1, \quad w_2 = 0.2, \quad w_3 = 0.5 \]

\section{Part (a): Sugeno (Zero-Order) Output}
In a zero-order Sugeno system, the consequents are singletons (constants):
\[ k_1 = 20, \quad k_2 = 50, \quad k_3 = 80 \]

The crisp output is calculated using the \textbf{Weighted Average (WA)} formula:

\[ z^*_{\text{Sugeno}} = \frac{\sum_{i=1}^{3} w_i k_i}{\sum_{i=1}^{3} w_i} \]

\textbf{Step 1: Calculate Numerator}
\[ \sum w_i k_i = (0.1 \times 20) + (0.2 \times 50) + (0.5 \times 80) = 2 + 10 + 40 = 52 \]

\textbf{Step 2: Calculate Denominator}
\[ \sum w_i = 0.1 + 0.2 + 0.5 = 0.8 \]

\textbf{Step 3: Final Output}
\[ z^*_{\text{Sugeno}} = \frac{52}{0.8} = \mathbf{65} \]

\section{Part (b): Mamdani Output (Comparison)}
For the Mamdani comparison, we use fuzzy set consequents centered at the same values ($c_i$) and apply the \textbf{Height Defuzzification} approximation (Weighted Centroid of Centers):
\[ c_1 = 20, \quad c_2 = 50, \quad c_3 = 80 \]

\[ z^*_{\text{Mamdani}} = \frac{\sum_{i=1}^{3} w_i c_i}{\sum_{i=1}^{3} w_i} \]

Since the weights $w_i$ and the centers $c_i$ are identical to the Sugeno parameters:
\[ z^*_{\text{Mamdani}} = \frac{52}{0.8} = \mathbf{65} \]

\section{Part (c): Decision Making}
\textbf{Comparison:} The outputs are \textbf{equal} ($65$) because the Height Defuzzification method mathematically reduces the Mamdani centroid calculation to a weighted average of centers, which matches the Sugeno formula exactly in this case.

\textbf{When to choose which method?}

\begin{enumerate}
    \item \textbf{Select Mamdani if:}
    \begin{itemize}
        \item \textbf{Interpretability} is the priority. You need to model human expert knowledge where rules must be readable (e.g., "IF Temp is High THEN Fan is Fast").
        \item The application requires capturing complex, non-linear input-output relationships that are best described linguistically.
    \end{itemize}
    
    \item \textbf{Select Sugeno if:}
    \begin{itemize}
        \item \textbf{Computational Efficiency} is critical. The weighted average calculation is significantly faster than calculating the centroid of an arbitrary area.
        \item You are using \textbf{optimization algorithms} (like genetic algorithms or neural networks) to tune the system parameters automatically (e.g., in ANFIS).
        \item You need to integrate the system with standard PID control or linear control theory.
    \end{itemize}
\end{enumerate}

\end{document}